% Options for packages loaded elsewhere
% Options for packages loaded elsewhere
\PassOptionsToPackage{unicode}{hyperref}
\PassOptionsToPackage{hyphens}{url}
\PassOptionsToPackage{dvipsnames,svgnames,x11names}{xcolor}
%
\documentclass[
  12pt,
  letterpaper,
  DIV=11,
  numbers=noendperiod]{scrartcl}
\usepackage{xcolor}
\usepackage[lmargin=2.5cm,rmargin=2.5cm,tmargin=2cm,bmargin=2cm]{geometry}
\usepackage{amsmath,amssymb}
\setcounter{secnumdepth}{5}
\usepackage{iftex}
\ifPDFTeX
  \usepackage[T1]{fontenc}
  \usepackage[utf8]{inputenc}
  \usepackage{textcomp} % provide euro and other symbols
\else % if luatex or xetex
  \usepackage{unicode-math} % this also loads fontspec
  \defaultfontfeatures{Scale=MatchLowercase}
  \defaultfontfeatures[\rmfamily]{Ligatures=TeX,Scale=1}
\fi
\usepackage{lmodern}
\ifPDFTeX\else
  % xetex/luatex font selection
  \setmainfont[]{Crimson Text}
  \setsansfont[]{Crimson Text}
  \setmonofont[]{Crimson Text}
\fi
% Use upquote if available, for straight quotes in verbatim environments
\IfFileExists{upquote.sty}{\usepackage{upquote}}{}
\IfFileExists{microtype.sty}{% use microtype if available
  \usepackage[]{microtype}
  \UseMicrotypeSet[protrusion]{basicmath} % disable protrusion for tt fonts
}{}
\makeatletter
\@ifundefined{KOMAClassName}{% if non-KOMA class
  \IfFileExists{parskip.sty}{%
    \usepackage{parskip}
  }{% else
    \setlength{\parindent}{0pt}
    \setlength{\parskip}{6pt plus 2pt minus 1pt}}
}{% if KOMA class
  \KOMAoptions{parskip=half}}
\makeatother
% Make \paragraph and \subparagraph free-standing
\makeatletter
\ifx\paragraph\undefined\else
  \let\oldparagraph\paragraph
  \renewcommand{\paragraph}{
    \@ifstar
      \xxxParagraphStar
      \xxxParagraphNoStar
  }
  \newcommand{\xxxParagraphStar}[1]{\oldparagraph*{#1}\mbox{}}
  \newcommand{\xxxParagraphNoStar}[1]{\oldparagraph{#1}\mbox{}}
\fi
\ifx\subparagraph\undefined\else
  \let\oldsubparagraph\subparagraph
  \renewcommand{\subparagraph}{
    \@ifstar
      \xxxSubParagraphStar
      \xxxSubParagraphNoStar
  }
  \newcommand{\xxxSubParagraphStar}[1]{\oldsubparagraph*{#1}\mbox{}}
  \newcommand{\xxxSubParagraphNoStar}[1]{\oldsubparagraph{#1}\mbox{}}
\fi
\makeatother

\usepackage{color}
\usepackage{fancyvrb}
\newcommand{\VerbBar}{|}
\newcommand{\VERB}{\Verb[commandchars=\\\{\}]}
\DefineVerbatimEnvironment{Highlighting}{Verbatim}{commandchars=\\\{\}}
% Add ',fontsize=\small' for more characters per line
\usepackage{framed}
\definecolor{shadecolor}{RGB}{241,243,245}
\newenvironment{Shaded}{\begin{snugshade}}{\end{snugshade}}
\newcommand{\AlertTok}[1]{\textcolor[rgb]{0.68,0.00,0.00}{#1}}
\newcommand{\AnnotationTok}[1]{\textcolor[rgb]{0.37,0.37,0.37}{#1}}
\newcommand{\AttributeTok}[1]{\textcolor[rgb]{0.40,0.45,0.13}{#1}}
\newcommand{\BaseNTok}[1]{\textcolor[rgb]{0.68,0.00,0.00}{#1}}
\newcommand{\BuiltInTok}[1]{\textcolor[rgb]{0.00,0.23,0.31}{#1}}
\newcommand{\CharTok}[1]{\textcolor[rgb]{0.13,0.47,0.30}{#1}}
\newcommand{\CommentTok}[1]{\textcolor[rgb]{0.37,0.37,0.37}{#1}}
\newcommand{\CommentVarTok}[1]{\textcolor[rgb]{0.37,0.37,0.37}{\textit{#1}}}
\newcommand{\ConstantTok}[1]{\textcolor[rgb]{0.56,0.35,0.01}{#1}}
\newcommand{\ControlFlowTok}[1]{\textcolor[rgb]{0.00,0.23,0.31}{\textbf{#1}}}
\newcommand{\DataTypeTok}[1]{\textcolor[rgb]{0.68,0.00,0.00}{#1}}
\newcommand{\DecValTok}[1]{\textcolor[rgb]{0.68,0.00,0.00}{#1}}
\newcommand{\DocumentationTok}[1]{\textcolor[rgb]{0.37,0.37,0.37}{\textit{#1}}}
\newcommand{\ErrorTok}[1]{\textcolor[rgb]{0.68,0.00,0.00}{#1}}
\newcommand{\ExtensionTok}[1]{\textcolor[rgb]{0.00,0.23,0.31}{#1}}
\newcommand{\FloatTok}[1]{\textcolor[rgb]{0.68,0.00,0.00}{#1}}
\newcommand{\FunctionTok}[1]{\textcolor[rgb]{0.28,0.35,0.67}{#1}}
\newcommand{\ImportTok}[1]{\textcolor[rgb]{0.00,0.46,0.62}{#1}}
\newcommand{\InformationTok}[1]{\textcolor[rgb]{0.37,0.37,0.37}{#1}}
\newcommand{\KeywordTok}[1]{\textcolor[rgb]{0.00,0.23,0.31}{\textbf{#1}}}
\newcommand{\NormalTok}[1]{\textcolor[rgb]{0.00,0.23,0.31}{#1}}
\newcommand{\OperatorTok}[1]{\textcolor[rgb]{0.37,0.37,0.37}{#1}}
\newcommand{\OtherTok}[1]{\textcolor[rgb]{0.00,0.23,0.31}{#1}}
\newcommand{\PreprocessorTok}[1]{\textcolor[rgb]{0.68,0.00,0.00}{#1}}
\newcommand{\RegionMarkerTok}[1]{\textcolor[rgb]{0.00,0.23,0.31}{#1}}
\newcommand{\SpecialCharTok}[1]{\textcolor[rgb]{0.37,0.37,0.37}{#1}}
\newcommand{\SpecialStringTok}[1]{\textcolor[rgb]{0.13,0.47,0.30}{#1}}
\newcommand{\StringTok}[1]{\textcolor[rgb]{0.13,0.47,0.30}{#1}}
\newcommand{\VariableTok}[1]{\textcolor[rgb]{0.07,0.07,0.07}{#1}}
\newcommand{\VerbatimStringTok}[1]{\textcolor[rgb]{0.13,0.47,0.30}{#1}}
\newcommand{\WarningTok}[1]{\textcolor[rgb]{0.37,0.37,0.37}{\textit{#1}}}

\usepackage{longtable,booktabs,array}
\usepackage{calc} % for calculating minipage widths
% Correct order of tables after \paragraph or \subparagraph
\usepackage{etoolbox}
\makeatletter
\patchcmd\longtable{\par}{\if@noskipsec\mbox{}\fi\par}{}{}
\makeatother
% Allow footnotes in longtable head/foot
\IfFileExists{footnotehyper.sty}{\usepackage{footnotehyper}}{\usepackage{footnote}}
\makesavenoteenv{longtable}
\usepackage{graphicx}
\makeatletter
\newsavebox\pandoc@box
\newcommand*\pandocbounded[1]{% scales image to fit in text height/width
  \sbox\pandoc@box{#1}%
  \Gscale@div\@tempa{\textheight}{\dimexpr\ht\pandoc@box+\dp\pandoc@box\relax}%
  \Gscale@div\@tempb{\linewidth}{\wd\pandoc@box}%
  \ifdim\@tempb\p@<\@tempa\p@\let\@tempa\@tempb\fi% select the smaller of both
  \ifdim\@tempa\p@<\p@\scalebox{\@tempa}{\usebox\pandoc@box}%
  \else\usebox{\pandoc@box}%
  \fi%
}
% Set default figure placement to htbp
\def\fps@figure{htbp}
\makeatother





\setlength{\emergencystretch}{3em} % prevent overfull lines

\providecommand{\tightlist}{%
  \setlength{\itemsep}{0pt}\setlength{\parskip}{0pt}}



 


\usepackage[dvipsnames]{xcolor}
\usepackage{amsfonts}
\usepackage{graphicx}
\usepackage{amsmath}
\usepackage{booktabs}
\usepackage{lscape}
\usepackage{adjustbox}
\usepackage{makecell}
\usepackage{threeparttable}
\usepackage{caption}

\newcommand{\blandscape}{\begin{landscape}}
\newcommand{\elandscape}{\end{landscape}}
\newcommand{\E}{\mathbb{E}}
\newcommand{\V}{\mathbb{V}}
\newcommand{\Cov}{\text{Cov}}
\newcommand\indep{\perp\!\!\!\perp}
\newcommand\nindep{\not\!\perp\!\!\!\perp}
\newcommand{\1}[1]{\mathds{1}\left[#1\right]}
\newcommand{\X}{\mathbf{X}}
\newcommand{\Y}{\mathbf{Y}}
\usepackage{booktabs}
\usepackage{longtable}
\usepackage{array}
\usepackage{multirow}
\usepackage{wrapfig}
\usepackage{float}
\usepackage{colortbl}
\usepackage{pdflscape}
\usepackage{tabu}
\usepackage{threeparttable}
\usepackage{threeparttablex}
\usepackage[normalem]{ulem}
\usepackage{makecell}
\usepackage{xcolor}
\KOMAoption{captions}{tableheading}
\makeatletter
\@ifpackageloaded{caption}{}{\usepackage{caption}}
\AtBeginDocument{%
\ifdefined\contentsname
  \renewcommand*\contentsname{Table of contents}
\else
  \newcommand\contentsname{Table of contents}
\fi
\ifdefined\listfigurename
  \renewcommand*\listfigurename{List of Figures}
\else
  \newcommand\listfigurename{List of Figures}
\fi
\ifdefined\listtablename
  \renewcommand*\listtablename{List of Tables}
\else
  \newcommand\listtablename{List of Tables}
\fi
\ifdefined\figurename
  \renewcommand*\figurename{Figure}
\else
  \newcommand\figurename{Figure}
\fi
\ifdefined\tablename
  \renewcommand*\tablename{Table}
\else
  \newcommand\tablename{Table}
\fi
}
\@ifpackageloaded{float}{}{\usepackage{float}}
\floatstyle{ruled}
\@ifundefined{c@chapter}{\newfloat{codelisting}{h}{lop}}{\newfloat{codelisting}{h}{lop}[chapter]}
\floatname{codelisting}{Listing}
\newcommand*\listoflistings{\listof{codelisting}{List of Listings}}
\makeatother
\makeatletter
\makeatother
\makeatletter
\@ifpackageloaded{caption}{}{\usepackage{caption}}
\@ifpackageloaded{subcaption}{}{\usepackage{subcaption}}
\makeatother
\usepackage{bookmark}
\IfFileExists{xurl.sty}{\usepackage{xurl}}{} % add URL line breaks if available
\urlstyle{same}
\hypersetup{
  pdftitle={Quantitative Methods -- Human Sciences},
  pdfauthor={Pablo Geraldo Bastías (pablo.geraldo@sociology.ox.ac.uk); Your name here},
  colorlinks=true,
  linkcolor={blue},
  filecolor={Maroon},
  citecolor={Blue},
  urlcolor={Blue},
  pdfcreator={LaTeX via pandoc}}


\title{Quantitative Methods -- Human Sciences}
\usepackage{etoolbox}
\makeatletter
\providecommand{\subtitle}[1]{% add subtitle to \maketitle
  \apptocmd{\@title}{\par {\large #1 \par}}{}{}
}
\makeatother
\subtitle{Difference-in-differences}
\author{Pablo Geraldo Bastías
(pablo.geraldo@sociology.ox.ac.uk) \and Your name here}
\date{}
\begin{document}
\maketitle


\newcommand{\benum}{\begin{enumerate}}
\newcommand{\eenum}{\end{enumerate}}
\newcommand{\bitem}{\begin{itemize}}
\newcommand{\eitem}{\end{itemize}}

\section{Difference-in-difference in theory (for
self-study)}\label{difference-in-difference-in-theory-for-self-study}

For this section, I ask you to explain the difference-in-difference
estimator for the Average Treatment Effect on the Treated (ATT)

\subsection{Write down and explain the parallel trends
assumptions}\label{write-down-and-explain-the-parallel-trends-assumptions}

\subsection{Is the parallel trend assumption verifiable from data? What
about ``pre-trend''
plots?}\label{is-the-parallel-trend-assumption-verifiable-from-data-what-about-pre-trend-plots}

\subsection{Is the two-way fixed effects (TWFE) estimator always
equivalent to the difference-in-differences (DiD)
model?}\label{is-the-two-way-fixed-effects-twfe-estimator-always-equivalent-to-the-difference-in-differences-did-model}

\section{DiD in practice}\label{did-in-practice}

In this problem, you will analyze data from the Card and Krueger paper
discussed in lecture.

\begin{quote}
Card, D. and A. B. Krueger (1994), "Minimum Wages and
Employment: A Case Study of the Fast-Food Industry in New Jersey and
Pennsylvania," \textit{American Economic Review}, vol. 84, 772-793.
\end{quote}

Card and Krueger are interested in estimating the impact of minimum wage
on teenage employment. Conventional economic wisdom states that raises
in minimum wages hurt employment, especially teenage employment, which
often takes wages that will be affected by minimum wage law. However,
empirical analysis has failed to find evidence of employment responses
to raises in minimum wages. In 1992, New Jersey's minimum wage increased
from \$4.25 to \$5.05 while the minimum wage in Pennsylvania remained at
\$4.25. The authors used data on employment at fast-food establishments
in New Jersey and Pennsylvania before and after the increase in the
minimum wage to measure the impact of the increase in minimum wage on
teenage employment. Download and import the dataset,
\texttt{card\_krueger.RData}, from the course website.

The dataset contains the following variables, with one row corresponding
to a single fast-food restaurant:

\item

\texttt{state}: New Jersey or Pennsylvania

\item

\texttt{chain}: the fast-food chain to which the restaurant belongs

\item

\texttt{wage\_pre}: starting wage in February 1992, in dollars per hour

\item

\texttt{wage\_post}: starting wage in November 1992, in dollars per hour

\item

\texttt{emp\_pre}: employment in February 1992, in number of full-time
equivalent (FTE) employees

\item

\texttt{emp\_post}: employment in November 1992, in number of full-time
equivalent (FTE) employees

\item

\texttt{closed}: whether the store was closed in November 1992

\begin{enumerate}
\item

Assume that the fast-food restaurants surveyed by Card and Krueger
represent a random sample from a larger population of all fast-food
restaurants in New Jersey and eastern Pennsylvania. Consider the
estimands in table \ref{tab:estimands}, which correspond to the mean
level of full-time equivalent (FTE) employment for population subgroups
(restaurants within a given state-time). For example,
\(\alpha = \mathbb{E}[Y_i|T=\text{February},\text{state}=\text{NJ}]\).

\begin{table}[!ht]
\begin{center}
\begin{tabular}{rcc}
               &   February    &   November    \\
New Jersey      &   $\alpha$      &   $\beta$       \\
Pennsylvania    &   $\gamma$      &   $\delta$      \\
\end{tabular}
\caption{Average population FTE employment, by state and time}
\label{tab:estimands}
\end{center}
\end{table}

\begin{enumerate}
\item[(i)]

Define the difference-in-differences (DID) estimand in terms of these
values, and concisely explain what this represents.

The DiD estimand can be obtained in two equivalent ways: 1) First
differencing post- and pre-treatment periods within states, and then
differencing states; 2) First differencing states within either pre- and
post-treatment periods, and then differencing both periods. This would
be equivalent to:

\[\text{DiD}_{(1)} = (\beta - \alpha) - (\delta - \gamma)\]
\[\text{DiD}_{(2)} = (\beta - \delta) - (\alpha - \gamma)\]

As the name indicates, this quantity is the combination of two
comparisons: over time, and over groups. A comparison over time in the
treatment group would be unbiased for the causal effect of interest (see
below) only if there is no time trend in its potential outcomes, while
an over groups comparison would be unbiased only if there is no
selection into the treatment (randomization). Both assumptions are hard
to believe in observational contexts, so by combining them we are trying
to weakening the identification assumption. In the context of DiD, the
result would be unbiased even in the presence of selection and time
trends, as long as 1) the time trend in the control condition is common
to both groups, and 2) the selection into the treatment (average gap in
control potential outcomes) is constant over time.

\item[(ii)]

Consider the eight potential quantities
\(\{\alpha(0), \alpha(1), \beta(0)...\}\). Let these represent the mean
\textit{potential} level of FTE employment levels that
\textit{would have realized} if the minimum wage had been raised in each
state at each time. For example,
\(\alpha(0) = \mathbb{E}[Y_i(D=0)|T=\text{February},\text{state}=\text{NJ}]\).
Define the causal quantity of interest, the ATT, in terms of these
potential outcomes. Describe which of these are observed.

The causal quantity of interest, the ATT, is the effect of rising the
minimun wage in the state that actually experienced the increase (New
Jersery). Then, our quantity of interest would be:

\[\E[Y_i(D=1)~|~T=\text{November}, \text{state}=\text{NJ}] - \E[Y_i(D=0)~|~T=\text{November}, \text{state}=\text{NJ}]\]

\[\beta(1) - \beta(0)\]

By consistency and SUTVA, we can replace \(\beta(1)\) by just \(\beta\),
the observed outcome of the treated group in the postreatment period.
The second term, the control potential outcome of the treated state in
the post-treatment period, is not observed (counterfactual).

\item[(iii)]

Discuss the identification assumption required for DID to yield a causal
effect: (a) state it in general terms and then as applied to this
problem; (b) explain in simple and intuitive terms why this is a useful
assumption; (c) finally, tell us whether you think it is a good
assumption here, and if not, what violations you imagine are
problematic.

The identification assumption required to give a causal interpretation
to DiD is known as parallel trends, and can be formally stated as:

\[\E[Y_0(1)-Y_0(0)~|~D=1] = \E[Y_0(1)-Y_0(0)~|~D=0]\]

The DiD is a combination of two comparisons: over time, and over groups.
A comparison over time in the treatment group would be unbiased for the
causal effect of interest only if there is no time trend in its
potential outcomes, while an over groups comparison would be unbiased
only if there is no selection into the treatment (randomization). Both
assumptions are hard to believe in observational contexts, so by
combining them we are are to weakening the identification assumption. In
the context of DiD, the result would be unbiased even in the presence of
selection and time trends, as long as 1) the time trend in the control
condition is common to both groups, and 2) the selection into the
treatment (average gap in control potential outcomes) is constant over
time.

The usefulness of this assumption in observational studies came from the
fact that one can accept some level of selection bias (in the sense of
non-interchangeable potential outcomes between groups), if that
selection is constant over time. In many contexts this assumption would
be more credible than assuming that, conditioning on certain covariates,
the treatment assignment is as-if random. Of course, this is still an
identification \textit{assumption}, in the sense that includes
unobserved or counterfactual quantities and, therefore, is non-testable.
We can only provide suggestive evidence of its plausibility and test
some observable implications, but not actually demonstrate that the
assumption holds.

In the context of the problem at hand, the assumption means that the
average change on employment between February and November in New
Jersey, given that there is no change in the minimum wage (control
condition), would be the same as the observed average change on
employment between February and November in Pennsylvania, where no
change in the minimum wage occurred during the period. Although this
parallel trends assumption is weaker than selection on observables, one
can imagine many circumstances that threaten its validity.

For example, potential violations would be that the trends in employment
in both states were significatively different to begin with, given
idiosyncratic reasons; that the closeness of both states makes that
policy change in one state affects the outcomes in the other state; that
other policies were changing at the same time that the rise in minimum
salary in New Jersery, and therefore that there are time-varying
confounders invalidating the assumptions, etc. Not being an expert in
the matter, I don't really know if the assumption is plausible enough,
but at least its plausibility should be explicitly defended because of
many potential invalidating situations.

\item[(iv)]

Making the required identifical assumption from above, show that the DID
estimand (part (i)) can identify the ATT (part (ii)). Show all your work
and note where the identification assumption comes in.

With a slightly different notation, we can write the DiD estimand as:

\begin{tabular}{l l l}
DiD = & $\E[Y(1)-Y(0)~|~D=1] - \E[Y(1)-Y(0)~|~D=0]$ & by definition\\
      & $\{\E[Y(1)~|~D=1]-\E[Y(0)~|~D=1]\} - \{\E[Y(1)~|~D=0]-\E[Y(0)~|~D=0]\}$ & linearity \\
      & $\{\E[Y_1(1)~|~D=1]-\E[Y_0(0)~|~D=1]\} - \{\E[Y_0(1)~|~D=0]-\E[Y_0(0)~|~D=0]\}$ & consistency \\ 
      & $\{\E[Y_1(1)~|~D=1]-\E[Y_0(0)~|~D=1]\} - \{\E[Y_0(1)~|~D=1]-\E[Y_0(0)~|~D=1]\}$ & parallel trends \\ 
      &  $\{\E[Y_1(1)~|~D=1]-\E[Y_0(1)~|~D=1]\} -\E[Y_0(0)~|~D=1]+\E[Y_0(0)~|~D=1]$ & rearrange \\
      & $\E[Y_1(1)~|~D=1]-\E[Y_0(1)~|~D=1]$ & cancell out\\
      & ATT
\end{tabular}

\end{enumerate}

\item

Estimate all quantities in table \ref{tab:estimands} and write them in a
similar table.

\begin{Shaded}
\begin{Highlighting}[]
\FunctionTok{library}\NormalTok{(}\StringTok{"kableExtra"}\NormalTok{)}
\FunctionTok{library}\NormalTok{(}\StringTok{"dplyr"}\NormalTok{)}
\end{Highlighting}
\end{Shaded}

\begin{verbatim}

Attaching package: 'dplyr'
\end{verbatim}

\begin{verbatim}
The following object is masked from 'package:kableExtra':

    group_rows
\end{verbatim}

\begin{verbatim}
The following objects are masked from 'package:stats':

    filter, lag
\end{verbatim}

\begin{verbatim}
The following objects are masked from 'package:base':

    intersect, setdiff, setequal, union
\end{verbatim}

\begin{Shaded}
\begin{Highlighting}[]
\FunctionTok{load}\NormalTok{(}\StringTok{"card\_krueger.Rdata"}\NormalTok{)}

\NormalTok{d }\SpecialCharTok{|\textgreater{}}  \FunctionTok{group\_by}\NormalTok{(state) }\SpecialCharTok{|\textgreater{}} 
  \FunctionTok{summarise}\NormalTok{(}\AttributeTok{february =} \FunctionTok{mean}\NormalTok{(emp\_pre, }\AttributeTok{na.rm=}\ConstantTok{TRUE}\NormalTok{),}
            \AttributeTok{november =} \FunctionTok{mean}\NormalTok{(emp\_post, }\AttributeTok{na.rm=}\ConstantTok{TRUE}\NormalTok{)) }\SpecialCharTok{|\textgreater{}} \FunctionTok{kable}\NormalTok{()}
\end{Highlighting}
\end{Shaded}

\begin{longtable}[]{@{}lrr@{}}
\toprule\noalign{}
state & february & november \\
\midrule\noalign{}
\endhead
\bottomrule\noalign{}
\endlastfoot
nj & 20.43941 & 21.02743 \\
pa & 23.33117 & 21.16558 \\
\end{longtable}

\item

Calculate your DID estimate (with a standard error and p-value) and
interpret the result.

\begin{Shaded}
\begin{Highlighting}[]
\FunctionTok{library}\NormalTok{(}\StringTok{"estimatr"}\NormalTok{)}
\FunctionTok{library}\NormalTok{(}\StringTok{"tidyr"}\NormalTok{)}

\CommentTok{\#d\_long \textless{}{-} gather(d, period, employment, emp\_pre:emp\_post) |\textgreater{}}
\CommentTok{\#  mutate(period = if\_else(period=="emp\_post", 1, 0),}
\CommentTok{\#         state = if\_else(state=="nj", 1, 0))}

\NormalTok{d\_long }\OtherTok{\textless{}{-}}\NormalTok{ d }\SpecialCharTok{|\textgreater{}}
  \FunctionTok{pivot\_longer}\NormalTok{(}
    \AttributeTok{cols =}\NormalTok{ emp\_pre}\SpecialCharTok{:}\NormalTok{emp\_post,}
    \AttributeTok{names\_to =} \StringTok{"period"}\NormalTok{,}
    \AttributeTok{values\_to =} \StringTok{"employment"}
\NormalTok{  ) }\SpecialCharTok{|\textgreater{}}
  \FunctionTok{mutate}\NormalTok{(}\AttributeTok{period =} \FunctionTok{if\_else}\NormalTok{(period }\SpecialCharTok{==} \StringTok{"emp\_post"}\NormalTok{, }\DecValTok{1}\NormalTok{, }\DecValTok{0}\NormalTok{),}
         \AttributeTok{state  =} \FunctionTok{if\_else}\NormalTok{(state }\SpecialCharTok{==} \StringTok{"nj"}\NormalTok{, }\DecValTok{1}\NormalTok{, }\DecValTok{0}\NormalTok{))}
\NormalTok{did }\OtherTok{\textless{}{-}}
  \FunctionTok{lm\_robust}\NormalTok{(employment }\SpecialCharTok{\textasciitilde{}}\NormalTok{ period }\SpecialCharTok{+}\NormalTok{ state }\SpecialCharTok{+}\NormalTok{ state}\SpecialCharTok{*}\NormalTok{period, }\AttributeTok{data=}\NormalTok{d\_long) }\SpecialCharTok{|\textgreater{}}
  \FunctionTok{print}\NormalTok{() }\SpecialCharTok{|\textgreater{}} \FunctionTok{data.frame}\NormalTok{() }
\end{Highlighting}
\end{Shaded}

\begin{Shaded}
\begin{Highlighting}[]
\NormalTok{did }\SpecialCharTok{|\textgreater{}} \FunctionTok{kable}\NormalTok{()}
\end{Highlighting}
\end{Shaded}

\begin{longtable}[]{@{}
  >{\raggedright\arraybackslash}p{(\linewidth - 14\tabcolsep) * \real{0.1625}}
  >{\raggedleft\arraybackslash}p{(\linewidth - 14\tabcolsep) * \real{0.1250}}
  >{\raggedleft\arraybackslash}p{(\linewidth - 14\tabcolsep) * \real{0.1375}}
  >{\raggedleft\arraybackslash}p{(\linewidth - 14\tabcolsep) * \real{0.1250}}
  >{\raggedleft\arraybackslash}p{(\linewidth - 14\tabcolsep) * \real{0.1250}}
  >{\raggedleft\arraybackslash}p{(\linewidth - 14\tabcolsep) * \real{0.1375}}
  >{\raggedleft\arraybackslash}p{(\linewidth - 14\tabcolsep) * \real{0.1375}}
  >{\raggedleft\arraybackslash}p{(\linewidth - 14\tabcolsep) * \real{0.0500}}@{}}
\toprule\noalign{}
\begin{minipage}[b]{\linewidth}\raggedright
\end{minipage} & \begin{minipage}[b]{\linewidth}\raggedleft
Estimate
\end{minipage} & \begin{minipage}[b]{\linewidth}\raggedleft
Std..Error
\end{minipage} & \begin{minipage}[b]{\linewidth}\raggedleft
t.value
\end{minipage} & \begin{minipage}[b]{\linewidth}\raggedleft
Pr\ldots t..
\end{minipage} & \begin{minipage}[b]{\linewidth}\raggedleft
CI.Lower
\end{minipage} & \begin{minipage}[b]{\linewidth}\raggedleft
CI.Upper
\end{minipage} & \begin{minipage}[b]{\linewidth}\raggedleft
DF
\end{minipage} \\
\midrule\noalign{}
\endhead
\bottomrule\noalign{}
\endlastfoot
(Intercept) & 23.331169 & 1.351149 & 17.267652 & 0.0000000 & 20.6789023
& 25.9834354 & 790 \\
period & -2.165584 & 1.647808 & -1.314222 & 0.1891530 & -5.4001833 &
1.0690145 & 790 \\
state & -2.891761 & 1.443583 & -2.003183 & 0.0454988 & -5.7254731 &
-0.0580484 & 790 \\
period:state & 2.753606 & 1.801200 & 1.528762 & 0.1267237 & -0.7820976 &
6.2893092 & 790 \\
\end{longtable}

Here our focal result is the third coefficient, \texttt{period:state},
which corresponds to the difference-in-difference. If our identification
assumption holds, then we can interpret the coefficient causally, as the
effect of increasing the minimum wage during 1992 in the state of New
Jersey, from 4.25 USD to 5.05 USD, which produces an average increase of
2.75 full-time equivalent employees in the fast-food chains affected by
the change. The result is not statistically significant at conventional
levels.

\end{enumerate}




\end{document}
